\hypertarget{appendix-24475e1596a8}{}\label{referencestart:appendix-24475e1596a8}
\section{Boolean algebra}\label{appendix-24475e1596a8-boolean-algebra}

\hypertarget{appendix-24475e1596a8-boolean-algebra}{}

\textbf{Curriculum depth:} Developed · \textbf{First seen:} C001 · \textbf{Latest development:} C002

\begin{allmightroledefinition}

\subsection{Definition and mechanism}\label{appendix-24475e1596a8-definition-and-mechanism}

\hypertarget{appendix-24475e1596a8-definition-and-mechanism}{}

Operations over false/true values can be defined and exhaustively checked with truth tables.

\end{allmightroledefinition}

The logical domain is false/true. A truth table enumerates all assignments, so two Boolean inputs need four rows and three need eight. AND requires both true, OR accepts either, XOR requires different inputs, and NOT reverses truth. These rules also apply independently to bit lanes.

\subsection{Prerequisites}\label{appendix-24475e1596a8-prerequisites}

\hypertarget{appendix-24475e1596a8-prerequisites}{}

\begin{itemize}
\tightlist
\item
  \hyperref[appendix-ddbb69998a2e]{Bit} --- A bit has two possible logical states.
\end{itemize}

\subsection{Used by}\label{appendix-24475e1596a8-used-by}

\hypertarget{appendix-24475e1596a8-used-by}{}

\begin{itemize}
\item
  \hyperref[appendix-b572b631069b]{Guard round and sticky bits} --- uses boolean algebra to establish this mechanism's representation or policy.
\item
  \hyperref[appendix-68b52453737c]{Bitwise AND} --- uses boolean algebra as part of its representation or reasoning.
\item
  \hyperref[appendix-a427ce8ce8a6]{Bitwise OR} --- uses boolean algebra as part of its representation or reasoning.
\item
  \hyperref[appendix-d49161831368]{Bitwise XOR} --- uses boolean algebra as part of its representation or reasoning.
\item
  \hyperref[appendix-8f38927cfc54]{Bitwise NOT} --- uses boolean algebra as part of its representation or reasoning.
\item
  \hyperref[appendix-fe6b457d68d8]{Logical versus bitwise operators} --- uses boolean algebra as part of its representation or reasoning.
\item
  \hyperref[appendix-9079d3659da4]{De Morgan\textquotesingle s laws} --- uses boolean algebra as part of its representation or reasoning.
\item
  \hyperref[appendix-d52d5f78ffd1]{Carry and borrow} --- uses boolean algebra as part of its representation or reasoning.
\item
  \hyperref[appendix-05db24ccc28a]{Exhaustive finite-state testing} --- uses boolean algebra as part of its representation or reasoning.
\end{itemize}

\subsection{Recall and next depth}\label{appendix-24475e1596a8-recall-and-next-depth}

\hypertarget{appendix-24475e1596a8-recall-and-next-depth}{}

Explain boolean algebra without the label, work the example above, then state the assumption whose removal would change the result.

\textbf{Pending:} predicates, normal forms, branchless evaluation, and digital circuits.

\subsection{Source and navigation}\label{appendix-24475e1596a8-source-and-navigation}

\hypertarget{appendix-24475e1596a8-source-and-navigation}{}

\hyperref[01-foundations-i-ch001-from-a-signal-to-an-interpretation]{C001 teaching chapter} · \hyperref[01-foundations-i-ch001-arithmetic-becomes-a-policy]{C002 development} · \hyperref[appendix-fd090d982a15]{Part definitions and mechanisms} · \href{../07-Vision-XR-and-Playable-Worlds/main.pdf\#nameddest=appendix-4ad47bdd269a}{Dependency Map (Vision, XR, and Playable Worlds, Reference notes: Dependency Map)}
